\section{Algebraic properties, tails, and variance}\label{sec:properties}

\subsection{Linearity and order}

\begin{theorem}[Nonnegativity]\label{thm:nonneg}
If \(X_i\ge0\) at every positive-weight outcome, then \(\E[X]\ge0\).
\end{theorem}

\begin{proof}
Every weighted summand \(w_iX_i\) is nonnegative.  Their finite sum is
nonnegative, and the denominator \(W\) is positive.  With rational codes,
clear the common positive denominator of the \(X_i\) and prove the resulting
natural inequality by repeated addition.
\end{proof}

\begin{theorem}[Additivity and homogeneity]\label{thm:linearity}
For rational observables \(X,Y\) and \(a\in\Q\),
\[
 \E[X+Y]=\E[X]+\E[Y],\qquad
 \E[aX]=a\E[X].
\]
\end{theorem}

\begin{proof}
For additivity,
\[
 W\E[X+Y]=\sum_iw_i(X_i+Y_i)
 =\sum_iw_iX_i+\sum_iw_iY_i
 =W(\E[X]+\E[Y]).
\]
The middle equality is the bounded-sum addition theorem proved by
quantifier-free induction.  Since \(W>0\), the rational codes on both sides are
equal.  Homogeneity follows by pulling the fixed factor \(a\) through each
finite addition.  Internally one first chooses a common denominator for
\(a,X_i,Y_i\), so only signed natural polynomial identities remain.
\end{proof}

\begin{corollary}[Finite linear combinations]\label{cor:finite-linearity}
For a coded list \((a_j,X_j)_{j<r}\),
\[
 \E\left[\sum_{j<r}a_jX_j\right]
 =\sum_{j<r}a_j\E[X_j].
\]
\end{corollary}

\begin{proof}
Use quantifier-free induction on \(r\).  The induction formula is an equality
of coded rationals.  The successor step is \Cref{thm:linearity}.
\end{proof}

\begin{corollary}[Expectation as a linear form]\label{cor:linear-form}
For a fixed finite law, \(X\mapsto\E[X]\) is a rational linear functional on
the finite rational module of observables.
\end{corollary}

\begin{theorem}[Monotonicity]\label{thm:monotone}
If \(X_i\le Y_i\) at every positive-weight outcome, then
\(\E[X]\le\E[Y]\).
\end{theorem}

\begin{proof}
The difference \(Y-X\) is pointwise nonnegative, so
\(0\le\E[Y-X]=\E[Y]-\E[X]\) by nonnegativity and linearity.
Equivalently, sum the pointwise inequalities directly.  The second proof is
slightly shorter in the object language because it avoids introducing a
signed subtraction term.
\end{proof}

\begin{corollary}[Constants and iterated expectation]\label{cor:constant}
For constant \(c\), \(\E[c]=c\).  Consequently the expectation of the constant
random variable \(\E[X]\) is again \(\E[X]\).
\end{corollary}

\begin{proof}
\(\sum_iw_ic=c\sum_iw_i=cW\).  Apply the same equation with
\(c=\E[X]\).
\end{proof}

\subsection{Positive and negative parts}

\begin{definition}[Parts and absolute value]\label{def:parts}
For a coded rational \(x\), define
\[
 x^+=\max(x,0),\qquad x^-=\max(-x,0),\qquad \abs{x}=x^++x^-.
\]
These operations are primitive recursive because rational order is decidable.
Apply them pointwise to observables.
\end{definition}

\begin{lemma}[Pointwise part identities]\label{lem:part-identities}
For every coded rational \(x\),
\[
 x=x^+-x^-,\qquad \abs{x}=x^++x^-,\qquad x^+x^-=0.
\]
\end{lemma}

\begin{proof}
Decide whether \(x\ge0\).  In the first branch, \(x^+=x\) and \(x^-=0\); in
the second, \(x^+=0\) and \(x^-=-x\).  This is a permitted intuitionistic
case split because the coded comparison is decidable.
\end{proof}

\begin{theorem}[Expectation and parts]\label{thm:expect-parts}
For every finite rational observable,
\[
 \E[X]=\E[X^+]-\E[X^-],\qquad
 \E[\abs X]=\E[X^+]+\E[X^-].
\]
Moreover,
\[
 \abs{\E[X]}\le\E[\abs X].
\]
\end{theorem}

\begin{proof}
The two equalities follow from \Cref{lem:part-identities} and linearity.  Since
\(-\abs X\le X\le\abs X\), monotonicity gives
\(-\E\abs X\le\E X\le\E\abs X\), which is equivalent to the final
inequality by decidable rational order.
\end{proof}

\begin{theorem}[Zero absolute expectation]\label{thm:zero-abs}
If \(\E[\abs X]=0\), then \(X=0\) almost surely.  More generally, if
\(\E[\abs{X-Y}]=0\), then \(X=Y\) almost surely.
\end{theorem}

\begin{proof}
After choosing a common denominator \(D\), the hypothesis is
\[
 \sum_{i<m}w_i\abs{a_i}=0,
\]
where \(X_i=a_i/D\).  Each term is a natural number.  Zero-sum elimination
proves \(w_i\abs{a_i}=0\) for every \(i\).  Thus every positive-weight outcome
has \(a_i=0\).  Equivalently, the weighted failure indicator has sum zero.
Apply the same argument to \(X-Y\).
\end{proof}

\begin{theorem}[Almost-sure congruence]\label{thm:ae-congruence}
If \(X=Y\) almost surely, then \(\E[X]=\E[Y]\).  If
\(g:\Q\to\Q\) is a fixed primitive-recursive map on the finite input codes,
then also \(g(X)=g(Y)\) almost surely and \(\E[g(X)]=\E[g(Y)]\).
\end{theorem}

\begin{proof}
At every outcome with positive weight the values agree, by zero-sum
elimination on the exceptional mass.  Hence the weighted summands agree and so
do their bounded sums.  Equality is preserved by applying the function code
\(g\), giving the second assertion.
\end{proof}

\subsection{Indicators and event algebra}

\begin{theorem}[Indicator expectation]\label{thm:indicator}
For every decidable finite event \(A\),
\[
 \E[\ind_A]=\Pp(A).
\]
For events \(A,B\),
\[
 \ind_{A\cup B}=\ind_A+\ind_B-\ind_{A\cap B},
\]
and therefore
\[
 \Pp(A\cup B)=\Pp(A)+\Pp(B)-\Pp(A\cap B).
\]
If \(A\cap B=\varnothing\), the intersection term vanishes.
\end{theorem}

\begin{proof}
The first formula is the definition.  For the second, decide the four possible
truth-value pairs of \(A(i),B(i)\) at each finite outcome.  The displayed
identity holds in each case.  Apply expectation and linearity.  The formal
version may avoid signed subtraction by moving the intersection term to the
left.
\end{proof}

\begin{corollary}[Finite union bound]\label{cor:union-bound}
For a coded list of events \(A_0,\ldots,A_{r-1}\),
\[
 \Pp\left(\bigcup_{j<r}A_j\right)\le\sum_{j<r}\Pp(A_j).
\]
\end{corollary}

\begin{proof}
Pointwise,
\(\ind_{\cup_jA_j}\le\sum_j\ind_{A_j}\).  Sum and use monotonicity and finite
linearity.  The union indicator is computed by bounded Boolean recursion.
\end{proof}

\subsection{Integer tail formulas}

\begin{theorem}[Nonnegative integer tail sum]\label{thm:tail-sum}
Let \(X_i\in\N\) and suppose \(X_i\le B\) for all outcomes.  Then
\[
 X_i=\sum_{k<B}\ind_{X_i>k}
\]
pointwise, and hence
\[
 \E[X]=\sum_{k<B}\Pp(X>k).
\]
Equivalently,
\(\E[X]=\sum_{k=1}^{B}\Pp(X\ge k)\).
\end{theorem}

\begin{proof}
For a fixed natural \(x\le B\), the indicators
\(\ind_{x>0},\ldots,\ind_{x>B-1}\) consist of exactly \(x\) ones followed by
zeros.  Prove
\(x=\sum_{k<B}\ind_{x>k}\) by quantifier-free induction on \(x\), or by
induction on \(B\) with the invariant truncated at \(\min(x,B)\).  Substitute
\(x=X_i\), multiply by \(w_i\), sum over \(i\), and exchange the two finite
sums.  The alternative form is the index shift \(k\mapsto k-1\).
\end{proof}

\begin{theorem}[Signed integer CDF formula]\label{thm:signed-tail}
Suppose \(-A\le X_i\le B\) for natural \(A,B\).  Then
\[
 \E[X]
 =\sum_{k=0}^{B-1}\Pp(X>k)
  -\sum_{k=0}^{A-1}\Pp(X\le -k-1).
\]
\end{theorem}

\begin{proof}
Apply \Cref{thm:tail-sum} to \(X^+\) and \(X^-\), and use
\(X=X^+-X^-\).  For natural \(k\),
\(X^->k\) is equivalent to \(-X>k\), hence to \(X\le-k-1\) because \(X\) is
integer-valued.  All comparisons are decidable.
\end{proof}

\begin{theorem}[Finite layer cake]\label{thm:finite-layer-cake}
If \(X_i\ge0\) are rational with common denominator \(D\) and
\(X_i\le B/D\), then
\[
 \E[X]=\frac1D\sum_{k<B}\Pp(DX>k).
\]
This is the exact finite precursor of
\(\E[X]=\int_0^\infty\Pp(X>t)\,dt\).
\end{theorem}

\begin{proof}
Apply \Cref{thm:tail-sum} to the natural-valued observable \(DX\), then divide
the resulting identity by \(D>0\).
\end{proof}

\subsection{Variance and covariance}

\begin{definition}[Variance and covariance]\label{def:variance}
Let \(\mu=\E[X]\) and \(\nu=\E[Y]\).  Define
\[
 \Var(X)=\E[(X-\mu)^2],\qquad
 \Cov(X,Y)=\E[(X-\mu)(Y-\nu)].
\]
\end{definition}

\begin{theorem}[Second-moment identities]\label{thm:variance-identity}
For finite rational observables,
\[
 \Var(X)=\E[X^2]-\E[X]^2,
\]
\[
 \Cov(X,Y)=\E[XY]-\E[X]\E[Y],
\]
and
\[
 \Var(X+Y)=\Var(X)+\Var(Y)+2\Cov(X,Y).
\]
\end{theorem}

\begin{proof}
Expand the relevant squares and products and use finite linearity and the
constant rule.  For example,
\begin{align*}
\E[(X-\mu)^2]
 &=\E[X^2-2\mu X+\mu^2]\\
 &=\E[X^2]-2\mu\E[X]+\mu^2
 =\E[X^2]-\mu^2.
\end{align*}
The formal proof is a polynomial identity after replacing each expectation by
its weighted numerator.  ProveIt's ring-normalization layer is designed for
exactly such denominator-cleared identities.
\end{proof}

\begin{theorem}[Variance is nonnegative and detects constants]\label{thm:variance-zero}
\(\Var(X)\ge0\).  Moreover \(\Var(X)=0\) if and only if \(X=\E[X]\) almost
surely.
\end{theorem}

\begin{proof}
The square \((X-\mu)^2\) is pointwise nonnegative.  If its expectation is zero,
zero-sum elimination shows the square is zero at every positive-weight
outcome, hence \(X_i=\mu\).  The converse is immediate.
\end{proof}

\begin{theorem}[The mean minimizes squared error]\label{thm:mean-minimizes}
For every rational constant \(a\),
\[
 \E[(X-a)^2]=\Var(X)+(\E[X]-a)^2.
\]
Consequently \(a=\E[X]\) minimizes the expected squared error; equality at
another \(a\) implies \(a=\E[X]\).
\end{theorem}

\begin{proof}
Write \(X-a=(X-\mu)+(\mu-a)\).  After squaring, the cross term has expectation
\(2(\mu-a)\E[X-\mu]=0\).  The remaining terms give the identity.  The final
claim follows from nonnegativity of the square \((\mu-a)^2\).
\end{proof}
